% !TeX root = ../Book4.tex
\begin{figure}[htbp]
	\centering
	\begin{tikzpicture}[font=\small,scale=1.05]
		\fill[GeometryBlueLight] (0,0) ellipse (3 and 1.65);
		\fill[white] (-1.2,0) circle (0.62);
		\fill[white] (1.2,0) circle (0.62);
		\draw[GeometryGray,thick] (0,0) ellipse (3 and 1.65);
		\draw[GeometryGray,thick] (-1.2,0) circle (0.62);
		\draw[GeometryGray,thick] (1.2,0) circle (0.62);
		\draw[GeometryBlue,very thick] (-3,0)--(-1.82,0);
		\draw[GeometryBlue,very thick] (-0.58,0)--(0.58,0);
		\draw[GeometryBlue,very thick] (1.82,0)--(3,0);
		\node[above] at (-2.4,0.05) {\(s_{13}\)};
		\node[above] at (0,0.05) {\(s_{12}\)};
		\node[above] at (2.4,0.05) {\(s_{23}\)};
		\node at (-1.2,0) {\(C_1\)};
		\node at (1.2,0) {\(C_2\)};
		\node[above] at (0,1.65) {\(C_3\)};
		\node at (0,1) {上侧六边形};
		\node at (0,-1) {下侧六边形};
		\node[below] at (-1.2,-0.62) {\(\ell_1/2\)};
		\node[below] at (1.2,-0.62) {\(\ell_2/2\)};
		\node[below] at (0,-1.65) {\(\ell_3/2\)};
	\end{tikzpicture}
	\caption[裤子的三条缝线与六边形分解]{三条蓝色缝线分别连接不同边界，把裤子分成上下两个直角六边形。每个六边形的三条边界弧长为 \(\ell_1/2,\ell_2/2,\ell_3/2\)，另三边为缝线。图为平面拓扑示意，所画圆周和长度比例不表示双曲测量。}
	\label{g4:fig:0604:pants-seams}
\end{figure}
